\chapter*{Preface}
\addcontentsline{toc}{chapter}{Preface}

Ce livre est un compagnon d'apprentissage personnel. Son objectif n'est pas de
reproduire un cours existant, mais de transformer chaque notion rencontree en
un savoir durable : intuition, lecture des formules, demonstration, experience
numerique, limites et application a l'ingenierie.

Les chapitres sont rediges en francais, avec une attention particuliere aux
questions qui surgissent naturellement pendant l'apprentissage. Pourquoi une
densite n'est-elle pas une probabilite? Pourquoi $\pi$ apparait-il dans une
formule sans cercle visible? Pourquoi la moyenne minimise-t-elle les distances
carrees? Ces questions ne sont pas des digressions. Elles constituent souvent
le chemin le plus direct vers une comprehension profonde.

La version presente inaugure l'ouvrage avec la distribution gaussienne et la
detection d'anomalies. Le depot associe conserve le code, les tests, la
tracabilite des sources privees et le journal editorial. Les contenus tiers ne
sont pas republies; seules des explications originales et des metadonnees de
tracabilite figurent dans l'ouvrage.

\vspace{1em}
\begin{flushright}
Guillaume Harbonnier\\
Diderot ML, 2026
\end{flushright}

